\chapter{GitHub and working with others}
\label{ch:github}

\section{Remotes}\label{remotes}

GitHub keeps a copy of your repository online.

\begin{lstlisting}[language=bash]
git remote add origin https://github.com/you/project.git
git push -u origin main     # send your commits
git pull                    # fetch others' commits
git clone <address>         # copy a project to your computer
\end{lstlisting}

This run pushed to a test copy on the same computer rather than to GitHub, so its address is shortened. A push to GitHub also shows progress lines, and asks you to sign in the first time.

\textbf{What you should see}

\begin{lstlisting}[language=text]
$ git push -u origin main
To .../empty.git
 * [new branch]      main -> main
branch 'main' set up to track 'origin/main'.

$ git remote -v
origin	.../empty.git (fetch)
origin	.../empty.git (push)

$ git pull
From .../empty
   b889be1..4912949  main       -> origin/main
Updating b889be1..4912949
Fast-forward
 results.txt | 1 +
 1 file changed, 1 insertion(+)
\end{lstlisting}

After a real push, your project\textquotesingle s page on GitHub shows the files, with the latest commit message beside each one. This is my pamphlet project\textquotesingle s page, with the account name and picture hidden:

\begin{figure}
\centering
\includegraphics[width=\linewidth]{img/github-repository.png}
\caption{A repository page on GitHub, with the account name and picture hidden.}
\end{figure}

Before making anything public, check it for secrets and for unpublished work. Git keeps history, so deleting a file later doesn\textquotesingle t remove it from past commits.

\section{Issues and pull requests}\label{issues-and-pull-requests}

\begin{itemize}
\tightlist
\item
  \textbf{An issue} is a note about something to fix or discuss.
\item
  \textbf{A pull request} proposes a change: someone works on a branch, then asks for it to be merged. Others can review it line by line, and automatic checks run on it.
\end{itemize}

\section{Making a project welcoming}\label{making-a-project-welcoming}

{\def\LTcaptype{none} % do not increment counter
\begin{longtable}[]{@{}
  >{\raggedright\arraybackslash}p{(\linewidth - 2\tabcolsep) * \real{0.5000}}
  >{\raggedright\arraybackslash}p{(\linewidth - 2\tabcolsep) * \real{0.5000}}@{}}
\toprule\noalign{}
\begin{minipage}[b]{\linewidth}\raggedright
File
\end{minipage} & \begin{minipage}[b]{\linewidth}\raggedright
Purpose
\end{minipage} \\
\midrule\noalign{}
\endhead
\bottomrule\noalign{}
\endlastfoot
\texttt{README.md} & What this is, and how to use it \\
\texttt{LICENSE} & What others may do with it. Without one, they have no clear right to reuse it \\
\texttt{CONTRIBUTING.md} & How to help \\
\texttt{CITATION.cff} & How to cite it \\
\texttt{.github/ISSUE\_TEMPLATE/} & Forms that ask for the right details \\
\end{longtable}
}

Remember people who don\textquotesingle t use GitHub, and offer them a simple way to reach you. The pamphlet\textquotesingle s \href{https://ahmadammar.github.io/marxist-economics/contribute/}{help page}\footnote{\texttt{https://ahmadammar.github.io/marxist-economics/contribute/}} is one example.

\section{Sites and releases}\label{sites-and-releases}

GitHub Pages publishes a website from a repository for free. A GitHub Release attaches files to a tag, while \href{https://zenodo.org}{Zenodo}\footnote{\texttt{https://zenodo.org}} can archive a release and give it a citable DOI.

